\section{Road Reparation}
\label{sec:cses-road-reparation}

\problemstatement{Given $n$ cities and $m$ possible roads, each with a
repair cost, choose roads of minimum total cost so that every city is
reachable from every other. Output the minimum cost, or \texttt{IMPOSSIBLE}
if the cities cannot all be connected.}

\begin{patternbox}
``Connect all cities at minimum total cost'' is a \textbf{minimum spanning
tree}. \textbf{Kruskal's algorithm} sorts roads by cost and adds each that
joins two different components, using Disjoint Set for the cycle test; if
fewer than $n-1$ roads are added, the graph is disconnected.
\end{patternbox}

\subsection*{Kruskal with Union-Find}

Sort all roads ascending by cost. Process each: if its endpoints are in
different components (checked via Disjoint Set), add it to the MST, union
the components, and accumulate the cost. Skip roads whose endpoints are
already connected (they would form a cycle). After processing, if exactly
$n-1$ roads were added, output the total; otherwise the graph could not be
fully connected, so print \texttt{IMPOSSIBLE}.

\begin{ideabox}[Cheapest Safe Edge First]
The MST is built greedily: the cheapest edge crossing between two
components is always safe to include. Disjoint Set makes ``same component?''
near-constant, and reaching $n-1$ edges certifies a spanning tree exists.
\end{ideabox}

\subsection*{Algorithm}

\begin{enumerate}
  \item Sort roads by cost. Initialise Disjoint Set.
  \item For each road, if it joins two components, union and add its cost;
        count added roads.
  \item If $n-1$ roads were added, output the total cost; else
        \texttt{IMPOSSIBLE}.
\end{enumerate}

\subsection*{Dry run}

\dryrun{$n=3$; roads $1\!-\!2(3)$, $2\!-\!3(5)$, $1\!-\!3(4)$.}

\begin{itemize}
  \item Sorted: $3,4,5$. Add $1\!-\!2(3)$; add $1\!-\!3(4)$; skip
        $2\!-\!3$ (cycle).
  \item Total $=3+4=7$, two roads $=n-1$: connected.
\end{itemize}

\subsection*{C++ implementation}

\begin{lstlisting}
#include <bits/stdc++.h>
using namespace std;

struct DSU {
    vector<int> parent, sz;
    DSU(int n) : parent(n + 1), sz(n + 1, 1) { iota(parent.begin(), parent.end(), 0); }
    int find(int x) { return parent[x] == x ? x : parent[x] = find(parent[x]); }
    bool unite(int a, int b) {
        a = find(a); b = find(b);
        if (a == b) return false;
        if (sz[a] < sz[b]) swap(a, b);
        parent[b] = a; sz[a] += sz[b];
        return true;
    }
};

int main() {
    int n, m;
    cin >> n >> m;
    vector<tuple<long long,int,int>> edges(m);
    for (auto& [w, a, b] : edges) cin >> a >> b >> w;
    sort(edges.begin(), edges.end());

    DSU dsu(n);
    long long total = 0;
    int used = 0;
    for (auto& [w, a, b] : edges)
        if (dsu.unite(a, b)) { total += w; used++; }

    if (used == n - 1) cout << total << "\n";
    else cout << "IMPOSSIBLE\n";
    return 0;
}
\end{lstlisting}

\begin{complexitybox}
\textbf{Time Complexity.} $O(m\log m)$ -- sorting dominates; unions are
near-constant. \textbf{Space Complexity.} $O(n+m)$.
\end{complexitybox}

\begin{takeawaybox}
Minimum-cost full connection is an MST; Kruskal plus Disjoint Set builds it
and, by counting added edges, detects disconnection. This is the canonical
network-design primitive.
\end{takeawaybox}
